% Fuente única de clase: Recursos/2025-10-09.pdf, páginas 1--10.
En la clase 15 estudiamos generación y condición noetheriana para módulos. Distinguimos las álgebras finitas de las de tipo finito, demostramos el teorema de la base de Hilbert y comenzamos el estudio de divisibilidad, asociados e irreducibles en dominios.

%%%%%%%%%%%%%%%%%%%%%%%%%%%%%%%%%%%%%%%%%%%%%%%%%%%%%%%%%%%%%%%%%%%%%%%%%%%%%%%
\section{Generadores y módulos noetherianos}
%%%%%%%%%%%%%%%%%%%%%%%%%%%%%%%%%%%%%%%%%%%%%%%%%%%%%%%%%%%%%%%%%%%%%%%%%%%%%%%
%%%%%%%%%%%%%%%%%%%%%%%%%%%%%%%%%%%%%%%%%%%%%%%%%%%%%%%%%%%%%%%%%%%%%%%%%%%%%%%
\subsection{El submódulo generado por un conjunto}
%%%%%%%%%%%%%%%%%%%%%%%%%%%%%%%%%%%%%%%%%%%%%%%%%%%%%%%%%%%%%%%%%%%%%%%%%%%%%%%
% P01.
Si $A\subseteq M$, la inclusión de conjuntos se extiende de manera única a un morfismo $\varphi_A:R^{\oplus A}\to M$. Su imagen es el \emph{submódulo generado por $A$}, denotado $\ang{A}$:
\[
\ang{A}=\set*{\sum_{a\in A}r_aa\given
 r_a\in R,\ r_a=0\text{ salvo para finitos }a}.
\]
Contiene $A$, porque $a=1_Ra$, y todo submódulo que contiene $A$ contiene todas esas combinaciones. Es, por tanto, el menor submódulo con esa propiedad. Para $A=\varnothing$ se obtiene cero.

\begin{defi}
Un módulo $M$ es \emph{finitamente generado} si $M=\ang{a_1,\ldots,a_n}$ para alguna familia finita. Equivalentemente, existe un homomorfismo sobreyectivo $R^n\to M$ para algún $n\geq0$.
\end{defi}
La equivalencia se sigue de la fórmula del mediador libre: $(r_1,\ldots,r_n)\mapsto\sum_i r_ia_i$. Recíprocamente, las imágenes de los vectores coordenados generan la imagen de todo morfismo desde $R^n$.

\begin{ejem}
El anillo $R=\mathbb Z[x_1,x_2,\ldots]$, con infinitas indeterminadas, es generado por $1$ como módulo sobre sí mismo. Cada polinomio sólo utiliza finitamente muchas variables. Su ideal $I=(x_1,x_2,\ldots)$ ofrece un contraste entre generación de un módulo y de sus submódulos.
\end{ejem}

\begin{ejer}[Un submódulo que no es finitamente generado]\label{ejer:C15-infinitas}
Demostrar que $I=(x_1,x_2,\ldots)$ no es finitamente generado como $R$-módulo. Es la comprobación III.6.14 citada en la página 1.
\end{ejer}
\begin{sol}
Supongamos que $I=(f_1,\ldots,f_s)$. Entre los $f_i$ aparecen sólo finitamente muchas variables; elegimos $x_j$ distinta de todas ellas. Cada $f_i\in I$ tiene término constante cero. Evaluemos todas las variables salvo $x_j$ en cero, conservando $x_j$ como indeterminada. Este homomorfismo $R\to\mathbb Z[x_j]$ envía cada $f_i$ a cero, de modo que anula el ideal que generan. Pero no anula $x_j$, que también debería pertenecer a ese ideal. Es una contradicción.
\end{sol}

\begin{defi}
Un módulo es \emph{noetheriano} si todos sus submódulos son finitamente generados.
\end{defi}
En particular, todo módulo noetheriano es finitamente generado, pues es submódulo de sí mismo. La vuelta falla por el ejercicio anterior. Si $R$ es conmutativo, sus submódulos como módulo regular son sus ideales, de modo que esta definición recupera la de anillo noetheriano de la clase 12.

%%%%%%%%%%%%%%%%%%%%%%%%%%%%%%%%%%%%%%%%%%%%%%%%%%%%%%%%%%%%%%%%%%%%%%%%%%%%%%%
\subsection{Submódulos, cocientes y sucesiones exactas}
%%%%%%%%%%%%%%%%%%%%%%%%%%%%%%%%%%%%%%%%%%%%%%%%%%%%%%%%%%%%%%%%%%%%%%%%%%%%%%%
% P02--P03 y P06.
\begin{prop}\label{prop:C15-extension}
Si $N\subseteq M$ es submódulo, $M$ es noetheriano si y sólo si lo son $N$ y $M/N$.
\end{prop}
\begin{proof}
Si $M$ es noetheriano, todo submódulo de $N$ también es submódulo de $M$ y es finitamente generado. Para un submódulo $\bar K\subseteq M/N$, la correspondencia proporciona $K\supseteq N$ con $\bar K=K/N$. Si $K=\ang{k_1,\ldots,k_t}$, sus clases generan $K/N$, pues
\[
x=\sum_i a_ik_i\quad\Longrightarrow\quad
x+N=\sum_i a_i(k_i+N).
\]
Así, $N$ y $M/N$ son noetherianos.

Para la vuelta, sea $K\subseteq M$ cualquier submódulo. Como $N$ es noetheriano, $K\cap N=\ang{\ell_1,\ldots,\ell_s}$. Por el segundo isomorfismo,
\[
K/(K\cap N)\cong(K+N)/N\subseteq M/N,
\]
de modo que ese cociente es finitamente generado. Elegimos representantes $k_1,\ldots,k_t\in K$ de un conjunto de generadores. Para $x\in K$ existen $r_i$ con $x-\sum_i r_ik_i\in K\cap N$. Entonces existen $s_j$ tales que
\[
x=\sum_i r_ik_i+\sum_j s_j\ell_j.
\]
Por tanto, $K=\ang{k_1,\ldots,k_t,\ell_1,\ldots,\ell_s}$. Como $K$ era arbitrario, $M$ es noetheriano.
\end{proof}
\begin{nota}[Conclusión de la vuelta]
La última línea manuscrita de la prueba en la página 3 concluye «$N$ es noetheriano». Eso ya era una hipótesis; el argumento prueba que \emph{$M$ es noetheriano}, porque todo submódulo $K$ de $M$ acaba de resultar finitamente generado.
\end{nota}

\begin{coro}\label{coro:C15-finito-noetheriano}
Sea $R$ un anillo conmutativo noetheriano. Todo módulo sobre $R$ que admita un conjunto finito de generadores es también noetheriano.
\end{coro}
\begin{proof}
El módulo regular $R$ es noetheriano. Por inducción y la proposición, también lo es $R^n$: el submódulo $R^{n-1}\oplus0$ tiene cociente $R$. El caso $n=0$ es el módulo cero. Si $R^n\twoheadrightarrow M$, el primer isomorfismo identifica $M$ con un cociente de $R^n$, que es noetheriano por la misma proposición.
\end{proof}

Una sucesión de módulos es \emph{exacta} en un término cuando la imagen de la flecha que llega coincide con el núcleo de la que sale. En una sucesión exacta corta
\[
0\longrightarrow N\xrightarrow{\varphi}M
\xrightarrow{\psi}P\longrightarrow0,
\]
$\varphi$ es inyectivo, $\psi$ es sobreyectivo y $\im\varphi=\ker\psi$. Por el primer isomorfismo, $N\cong\im\varphi$ y $P\cong M/\im\varphi$. Así, la \cref{prop:C15-extension} equivale a que $M$ sea noetheriano si y sólo si lo son los dos extremos $N,P$.
\begin{nota}[Identificación del núcleo en una sucesión]
En la página 6, la aplicación con dominio $M$ y destino $P$ es $\psi$, así que el cociente que da $P$ es $M/\ker\psi$, no $M/\ker\varphi$ como aparece en una línea manuscrita. La expresión posterior $M/N$ supone identificar $N$ con su imagen mediante $\varphi$.
\end{nota}

%%%%%%%%%%%%%%%%%%%%%%%%%%%%%%%%%%%%%%%%%%%%%%%%%%%%%%%%%%%%%%%%%%%%%%%%%%%%%%%
\subsection{Cadenas ascendentes y elementos maximales}
%%%%%%%%%%%%%%%%%%%%%%%%%%%%%%%%%%%%%%%%%%%%%%%%%%%%%%%%%%%%%%%%%%%%%%%%%%%%%%%
% P06--P08.
\begin{prop}\label{prop:C15-cadenas}
Para un módulo $M$ son equivalentes:
\begin{enumerate}[label=\textup{(\roman*)}]
\item $M$ es noetheriano.
\item Toda cadena ascendente $N_1\subseteq N_2\subseteq\cdots$ de submódulos se estabiliza: existe $k$ con $N_j=N_k$ para todo $j\geq k$.
\item Toda familia no vacía de submódulos tiene un elemento maximal respecto de la inclusión.
\end{enumerate}
\end{prop}
\begin{proof}
Para (i)$\Rightarrow$(ii), la unión $N=\bigcup_i N_i$ es submódulo. Dados dos elementos, ambos pertenecen a algún término común de la cadena, por lo que su diferencia también pertenece; la estabilidad bajo escalares se comprueba en el término que contiene al elemento. Como $M$ es noetheriano, $N=\ang{a_1,\ldots,a_s}$. Todos los generadores están en algún $N_k$, tomando el máximo de sus índices. Entonces $N\subseteq N_k\subseteq N$, y $N_j=N_k$ para $j\geq k$. Si $N=0$, la conclusión es inmediata.

Para (ii)$\Rightarrow$(iii), si una familia no vacía no tuviera maximal, podríamos elegir $N_1$ en ella y, sucesivamente, un miembro $N_{i+1}$ que contuviera estrictamente a $N_i$. La cadena nunca se estabilizaría, contradiciendo (ii).

Para (iii)$\Rightarrow$(i), fijemos un submódulo \emph{arbitrario} $K\subseteq M$. La familia de sus submódulos finitamente generados es no vacía, porque contiene cero. Sea $L=\ang{a_1,\ldots,a_s}$ un maximal en esa familia. Para cualquier $x\in K$, $L+Rx=\ang{a_1,\ldots,a_s,x}$ también pertenece a ella y contiene $L$, luego es $L$. Así $x\in L$, con lo que $K=L$ es finitamente generado. Esto vale para todo $K$ y prueba (i).
\end{proof}
\begin{nota}[Aplicar maximalidad en cada submódulo]
La página 8 aplica el último argumento a los submódulos finitamente generados de $M$ y obtiene que $M$ es finitamente generado. Para concluir que es \emph{noetheriano}, hay que repetirlo dentro de cada $K\subseteq M$, como se hizo arriba; demostrar sólo la generación finita de $M$ no basta, por el \cref{ejer:C15-infinitas}.
\end{nota}

%%%%%%%%%%%%%%%%%%%%%%%%%%%%%%%%%%%%%%%%%%%%%%%%%%%%%%%%%%%%%%%%%%%%%%%%%%%%%%%
\section{Álgebras finitas y teorema de Hilbert}
%%%%%%%%%%%%%%%%%%%%%%%%%%%%%%%%%%%%%%%%%%%%%%%%%%%%%%%%%%%%%%%%%%%%%%%%%%%%%%%
%%%%%%%%%%%%%%%%%%%%%%%%%%%%%%%%%%%%%%%%%%%%%%%%%%%%%%%%%%%%%%%%%%%%%%%%%%%%%%%
\subsection{Finita como módulo y de tipo finito como álgebra}
%%%%%%%%%%%%%%%%%%%%%%%%%%%%%%%%%%%%%%%%%%%%%%%%%%%%%%%%%%%%%%%%%%%%%%%%%%%%%%%
% P04.
En este apartado $R$ y la $R$-álgebra $S$ son conmutativos, con homomorfismo estructural $\alpha:R\to S$.
\begin{defi}
Se dice que $S$ es una \emph{$R$-álgebra finita} si es finitamente generada como $R$-módulo, es decir, si existe una sobreyección lineal $R^n\twoheadrightarrow S$. Es \emph{de tipo finito} si es finitamente generada como álgebra, es decir, si existe un homomorfismo sobreyectivo de $R$-álgebras $R[x_1,\ldots,x_n]\twoheadrightarrow S$.
\end{defi}
La palabra «finita» se refiere aquí a la generación del módulo, no a la cardinalidad del conjunto. Usamos álgebras conmutativas para que la presentación por los polinomios conmutativos tenga la propiedad universal de la clase 14.

\begin{prop}
Toda $R$-álgebra conmutativa finita es de tipo finito.
\end{prop}
\begin{proof}
Si $a_1,\ldots,a_n$ generan $S$ como módulo, la propiedad universal da el morfismo
\[
\varphi:R[x_1,\ldots,x_n]\longrightarrow S,\qquad\varphi(x_i)=a_i.
\]
Todo $s\in S$ se escribe $s=\sum_i\alpha(r_i)a_i=\varphi(\sum_i r_ix_i)$, así que $\varphi$ es sobreyectiva.
\end{proof}

\begin{ejem}
Si $R\ne0$, $R[x]$ es de tipo finito, generado como álgebra por $x$, pero no es finito como módulo. Si unos polinomios $f_1,\ldots,f_s$ lo generaran como módulo, habría una cota $d$ para sus grados. Toda combinación con coeficientes en $R$ tendría grado a lo sumo $d$, mientras que $x^{d+1}$ tiene coeficiente principal $1_R\ne0$. No se puede generar así. Si no hubiera generadores no nulos, ni siquiera se generaría $1_R$.
\end{ejem}
\begin{nota}[Excepción del anillo cero]
El ejemplo de la página 4 requiere $R\ne0$. Si $R=0$, también $R[x]=0$ y sí es finitamente generado como módulo. La condición se explicita para mantener la convención unitaria de estas notas.
\end{nota}

%%%%%%%%%%%%%%%%%%%%%%%%%%%%%%%%%%%%%%%%%%%%%%%%%%%%%%%%%%%%%%%%%%%%%%%%%%%%%%%
\subsection{El ideal de coeficientes principales}
%%%%%%%%%%%%%%%%%%%%%%%%%%%%%%%%%%%%%%%%%%%%%%%%%%%%%%%%%%%%%%%%%%%%%%%%%%%%%%%
% P08: comprobación V.1.6 citada en la prueba de Hilbert.
Volvemos a anillos conmutativos $R$ arbitrarios. Para un ideal $I\subseteq R[x]$, sea
\[
A=\{0\}\cup\set{a\in R\given
a\text{ es coeficiente principal de algún }f\in I,\ f\ne0}.
\]
\begin{ejer}[Coeficientes principales de un ideal]\label{ejer:C15-principales}
Probar que $A$ es un ideal de $R$, como propone el ejercicio V.1.6 citado en la página 8.
\end{ejer}
\begin{sol}
Contiene cero. Si $a,b\ne0$ son coeficientes principales de $f,g\in I$ de grados $m,n$, elegimos $d=\max(m,n)$. El polinomio $x^{d-m}f-x^{d-n}g$ pertenece a $I$ y su coeficiente de grado $d$ es $a-b$. Si éste no es cero, es su coeficiente principal; si es cero, ya pertenece a $A$ por definición. Cuando uno de $a,b$ es cero se usa $f$ o $-g$. Así, $A$ es subgrupo aditivo.

Para $r\in R$, si $ra\ne0$, el polinomio $rf$ tiene ese coeficiente principal y pertenece a $I$. Si $ra=0$, el resultado pertenece a $A$ sin necesidad de afirmar nada sobre el grado de $rf$. Por tanto, $rA\subseteq A$ y $A$ es ideal, aun cuando $R$ tenga divisores de cero.
\end{sol}

%%%%%%%%%%%%%%%%%%%%%%%%%%%%%%%%%%%%%%%%%%%%%%%%%%%%%%%%%%%%%%%%%%%%%%%%%%%%%%%
\subsection{Teorema de la base de Hilbert}
%%%%%%%%%%%%%%%%%%%%%%%%%%%%%%%%%%%%%%%%%%%%%%%%%%%%%%%%%%%%%%%%%%%%%%%%%%%%%%%
% P08--P09.
\begin{teo}[Teorema de la base de Hilbert]\label{teo:C15-Hilbert}
Si $R$ es un anillo conmutativo noetheriano, también lo es $R[x]$. En consecuencia, toda $R$-álgebra conmutativa de tipo finito es noetheriana.
\end{teo}
\begin{proof}
Sea $I$ un ideal de $R[x]$. Si $I=0$, ya es principal. En otro caso, el ideal $A$ del ejercicio anterior es no nulo y, por noetherianidad, admite generadores no nulos $a_1,\ldots,a_t$. Elegimos $f_i\in I$ con coeficiente principal $a_i$, ponemos $d_i=\deg f_i$ y $d=\max_i d_i$.

El módulo $L=\ang{1,x,\ldots,x^{d-1}}\subseteq R[x]$ consiste en los polinomios de grado menor que $d$. Es finitamente generado y, por el \cref{coro:C15-finito-noetheriano}, es noetheriano como $R$-módulo. Por tanto, $L\cap I$ tiene generadores $g_1,\ldots,g_s$ como módulo. Si $d=0$, $L=0$ y esta lista es vacía. Afirmamos, como en V.1.4 de la pizarra, que
\[
I=(f_1,\ldots,f_t,g_1,\ldots,g_s)
\]
como ideal de $R[x]$.

Una inclusión es inmediata, pues todos esos elementos están en $I$. Para la otra, tomemos $h\ne0$ en $I$. Si $e=\deg h\geq d$, su coeficiente principal $a$ pertenece a $A$, de modo que $a=\sum_i b_ia_i$, con $b_i\in R$. Como $e\geq d_i$, el polinomio
\[
h-\sum_i b_ix^{e-d_i}f_i
\]
pertenece a $I$ y tiene grado menor que $e$, o es cero: su coeficiente de grado $e$ es $a-\sum_i b_ia_i=0$. Iteramos. El descenso estricto del grado garantiza que terminamos con
\[
h-\sum_i\beta_i(x)f_i\in L\cap I.
\]
El resultado es una combinación $\sum_j c_jg_j$ con $c_j\in R$, por lo que $h$ pertenece al ideal afirmado. El mismo argumento cubre los $h$ de grado menor que $d$ desde el principio, y cero pertenece automáticamente. Así, todo ideal es finitamente generado.

Aplicamos la primera parte sucesivamente para obtener que $R[x_1,\ldots,x_n]$ es noetheriano. Si $S$ es de tipo finito, es isomorfo a un cociente de ese anillo. Los ideales de un cociente corresponden a los ideales que contienen al núcleo y sus generadores proyectan a generadores del ideal cociente. Por ello $S$ es noetheriano como anillo.
\end{proof}
La última afirmación no dice que $S$ sea noetheriano como $R$-módulo: por ejemplo, sobre un campo no nulo $K$, $K[x]$ es un anillo noetheriano, pero no es finitamente generado como $K$-módulo.

%%%%%%%%%%%%%%%%%%%%%%%%%%%%%%%%%%%%%%%%%%%%%%%%%%%%%%%%%%%%%%%%%%%%%%%%%%%%%%%
\section{Divisibilidad e irreducibilidad en dominios}
%%%%%%%%%%%%%%%%%%%%%%%%%%%%%%%%%%%%%%%%%%%%%%%%%%%%%%%%%%%%%%%%%%%%%%%%%%%%%%%
%%%%%%%%%%%%%%%%%%%%%%%%%%%%%%%%%%%%%%%%%%%%%%%%%%%%%%%%%%%%%%%%%%%%%%%%%%%%%%%
\subsection{Panorama de la factorización}
%%%%%%%%%%%%%%%%%%%%%%%%%%%%%%%%%%%%%%%%%%%%%%%%%%%%%%%%%%%%%%%%%%%%%%%%%%%%%%%
% P05: portada del capítulo V y esquema orientativo.
La página 5 abre el capítulo V del libro con un esquema de clases de dominios. Lo conservamos como orientación: los rótulos «con factorización», «DFU» y «euclidiano» se precisan en la clase 16. La figura anuncia las relaciones que motivan ese estudio; no sustituye sus demostraciones.
% Esquema orientativo de 2025-10-09.pdf P05 y 2025-10-14.pdf P02.
% Los óvalos DFU y noetherianos se intersectan, sin contenerse mutuamente.

\begin{tikzpicture}[font=\small,line width=.6pt]
\filldraw[fill=grisCIMAT!30] (0,0) ellipse (7cm and 4cm);
\filldraw[fill=grisCIMAT!15] (0,-.35) ellipse (6.5cm and 3.4cm);
\draw[fill=white,fill opacity=.35,draw opacity=1] (-1.2,-.85) ellipse (4.5cm and 2.8cm);
\draw[fill=white,fill opacity=.35,draw opacity=1] (1.2,-.85) ellipse (4.5cm and 2.8cm);
\filldraw[fill=grisCIMAT!10] (0,-1.6) ellipse (2.5cm and 1.8cm);
\filldraw[fill=white] (0,-2) ellipse (1.9cm and 1.2cm);
\filldraw[fill=grisCIMAT!8] (0,-2.5) ellipse (1.1cm and .55cm);
\node[font=\normalsize\itshape] at (0,3.35) {Dominios enteros};
\node[font=\normalsize\itshape] at (0,2.4) {Dominios con factorización};
\node[rotate=18] at (-3.65,1.1) {DFU};
\node[rotate=-18] at (3.3,1.08) {Dominios noetherianos};
\node at (0,-.35) {DIP};
\node[font=\footnotesize] at (0,-1.45) {Dominios euclidianos};
\node at (0,-2.5) {Campos};
\end{tikzpicture}


%%%%%%%%%%%%%%%%%%%%%%%%%%%%%%%%%%%%%%%%%%%%%%%%%%%%%%%%%%%%%%%%%%%%%%%%%%%%%%%
\subsection{Divisores, asociados y elementos primos}
%%%%%%%%%%%%%%%%%%%%%%%%%%%%%%%%%%%%%%%%%%%%%%%%%%%%%%%%%%%%%%%%%%%%%%%%%%%%%%%
% P10.
\begin{defi}
En un anillo conmutativo, $a$ \emph{divide} a $b$, escrito $a\mid b$, si existe $c\in R$ con $b=ac$, es decir, si $b\in(a)$. Dos elementos son \emph{asociados} si $(a)=(b)$, equivalentemente, si cada uno divide al otro.
\end{defi}
La divisibilidad es reflexiva y transitiva por las identidades $a=a1$ y $b=ac$, $d=be\Rightarrow d=a(ce)$. Además, $a\mid b$ equivale a $(b)\subseteq(a)$, de modo que la dirección de inclusión de ideales es opuesta a la divisibilidad.

\begin{lema}\label{lema:C15-asociados}
En un dominio, dos elementos no nulos son asociados si y sólo si uno es múltiplo del otro por una unidad.
\end{lema}
\begin{proof}
Si $b=ac$ y $a=bd$, entonces $a=acd$. Cancelando $a\ne0$, $cd=1$, así que $d$ es unidad y $a=db$. Recíprocamente, $a=ub$ da $a\in(b)$ y $b=u^{-1}a\in(a)$, por lo que los ideales coinciden.
\end{proof}

\begin{defi}
En un dominio $R$, un elemento $a$ es \emph{primo} si $(a)$ es un ideal primo. Equivalentemente, $a$ no es unidad y
\[
a\mid bc\quad\Longrightarrow\quad a\mid b\ \text{o}\ a\mid c.
\]
Es \emph{irreducible} si no es unidad y toda factorización $a=bc$ tiene al menos un factor unidad.
\end{defi}
La equivalencia para primos es la \cref{prop:C13-criterios}. Esta convención del fragmento permite que $0$ sea primo en un dominio, porque $(0)$ es primo. En cambio, $0$ no es irreducible: $0=0\cdot0$ y ninguno de esos factores es unidad en un dominio no nulo. Por eso se exigirá explícitamente «primo no nulo» al comparar ambas nociones.

\begin{ejer}[Irreducibilidad e ideales principales]\label{ejer:C15-irreducible}
Para un elemento no nulo y no unidad $a$ de un dominio, demostrar que las condiciones siguientes equivalen a su irreducibilidad:
\begin{enumerate}[label=\textup{(\roman*)}]
\item Si $a=bc$, entonces $a$ es asociado de $b$ o de $c$.
\item Si $a=bc$, entonces $(a)=(b)$ o $(a)=(c)$.
\item Si $(a)\subseteq(b)$, entonces $(b)=(a)$ o $(b)=R$.
\end{enumerate}
Así, $(a)$ es maximal entre los ideales principales propios. La tercera formulación es la comprobación V.1.12 citada en la página 10.
\end{ejer}
\begin{sol}
Si $a$ es irreducible y $a=bc$, uno de los factores es unidad. Si $c$ es unidad, $a$ es asociado de $b$; si $b$ es unidad, es asociado de $c$. Recíprocamente, si por ejemplo $a=ub$ con $u$ unidad y también $a=bc$, cancelar $b\ne0$ da $c=u$. Así, (i) implica irreducibilidad. La equivalencia entre (i) y (ii) es la definición de asociados.

Si $a$ es irreducible y $(a)\subseteq(b)$, escribimos $a=bc$. Si $c$ es unidad, $(a)=(b)$; si $b$ es unidad, $(b)=R$. Esto prueba (iii). Para la vuelta, de $a=bc$ obtenemos $(a)\subseteq(b)$. Si $(b)=R$, $b$ es unidad. Si $(b)=(a)$, existe $d$ con $b=ad$, y $a=adc$; cancelar $a\ne0$ da $dc=1$, por lo que $c$ es unidad. Finalmente, (iii) dice exactamente que ningún ideal principal propio contiene estrictamente a $(a)$.
\end{sol}

\begin{obs}[Orientación bibliográfica]
Se consultaron los fragmentos originales de \citet{aluffi2009}: III.6.4--5, ejemplo III.6.5, definición III.6.6, proposición III.6.7, corolario III.6.8 y ejemplo III.6.9; la apertura del capítulo V; la proposición V.1.1, el teorema V.1.2, el lema V.1.3 y la afirmación V.1.4; el apartado V.1.2, lema V.1.5 y definición V.1.6. Se resuelven las comprobaciones III.6.14, V.1.6 y V.1.12 remitidas por los fragmentos. El caso del ideal cero y el de grado máximo cero se incorporan a la prueba de Hilbert. Las referencias explícitas a III.4 y III.6 en estas páginas permiten además corregir los números de capítulo que se habían atribuido erróneamente en clases anteriores.
\end{obs}
